% 霍普夫代数（综述）
% license CCBYSA3
% type Wiki

本文根据 CC-BY-SA 协议转载翻译自维基百科\href{https://en.wikipedia.org/wiki/Hopf_algebra}{相关文章}

在数学中，霍普夫代数（Hopf algebra，以海因茨·霍普夫 Heinz Hopf 命名）是一种结构，它既是一个（含幺结合的）代数，又是一个（含余元的余结合）余代数；这些结构的相容性使它成为一个双代数，此外它还配备了一个满足特定性质的反同态。霍普夫代数的表示论尤其优美，因为兼容的余乘法、余元以及反元的存在，使得能够构造表示的张量积、平凡表示以及对偶表示。

霍普夫代数自然地出现在代数拓扑中（其起源之处，并与 H-空间概念相关）、群概形理论、群论（通过群环的概念），以及许多其他地方，因此可能是最常见的一类双代数。霍普夫代数也作为独立的研究对象受到关注，一方面有许多关于具体类实例的研究，另一方面也存在分类问题。它们拥有广泛的应用，从凝聚态物理与量子场论，到弦论以及大型强子对撞机（LHC）现象学。
\subsection{形式定义}
形式上，霍普夫代数是域$K$上的一个（结合且余结合的）双代数$H$，并配备一个$K$-线性映射$S: H \to H$（称为反元，antipode），使得下列图表交换：
\begin{figure}[ht]
\centering
\includegraphics[width=8cm]{./figures/080fad97a4c9a060.png}
\caption{} \label{fig_HPFds_1}
\end{figure}
这里，$\Delta$是双代数的余乘法，$\nabla$是其乘法，$\eta$ 是其单位元映射，$\varepsilon$是其余元映射。在无求和的Sweedler 记号下，这一性质也可以表述为：
$$
S(c_{(1)})c_{(2)} \;=\; c_{(1)}S(c_{(2)}) \;=\; \varepsilon(c)\,1 
\qquad \text{对于所有 } c \in H.~
$$
与代数的情况类似，在上述定义中，可以将底层的域$K$替换为交换环$R$。【4】

霍普夫代数的定义是自对偶的（这在上面图表的对称性中有所体现），因此如果可以定义 $H$的对偶（当$H$是有限维时总是可能的），那么该对偶也自动是一个霍普夫代数。【5】
\subsubsection{结构常数}
在固定了底层向量空间的一组基$\{ e_k \}$之后，可以通过乘法的结构常数来定义代数：
$$
e_i \nabla e_j = \sum_k \mu_{ij}^k e_k~
$$
对于余乘法：
$$
\Delta e_i = \sum_{j,k} \nu_i^{\,jk} \, e_j \otimes e_k~
$$
以及反元：
$$
S e_i = \sum_j \tau_i^{\,j} e_j~
$$
结合律要求：
$$
\mu_{ij}^k \, \mu_{kn}^m \;=\; \mu_{jn}^k \, \mu_{ik}^m~
$$
而余结合律要求：
$$
\nu_k^{\,ij} \, \nu_i^{\,mn} \;=\; \nu_k^{\,mi} \, \nu_i^{\,nj}~
$$
连接公理要求：
$$
\nu_k^{\,ij} \, \tau_j^{\,m} \, \mu_{im}^n 
\;=\; 
\nu_k^{\,jm} \, \tau_j^{\,i} \, \mu_{im}^n~
$$
\subsubsection{反元的性质}
反元$S$有时被要求具有一个$K$-线性逆映射，这在有限维情形下会自动成立【需要澄清】，或者当$H$是交换的、余交换的（或更一般地是拟三角的 quasitriangular）时也成立。

一般而言，$S$是一个反同态【6】，因此$S^2$是同态；若 $S$ 可逆（通常要求如此），那么 $S^2$ 就是一个自同构。

若$S^2 = \mathrm{id}_H$，则称该霍普夫代数为对合的，其底层带对合的代数是一个 $*$-代数。若$H$是有限维、在零特征域上半单、交换或余交换的，那么它是对合的。

如果一个双代数$B$承认一个反元$S$，那么这个$S$是唯一的（“一个双代数至多承认一个霍普夫代数结构”）【7】。因此，反元并不是可以随意选择的额外结构：作为霍普夫代数，是双代数的一个性质。

反元可以类比为群上的逆映射，即把元素$g$映射为$g^{-1}$【8】。
\subsubsection{霍普夫子代数}
若$H$是一个霍普夫代数，$A$是其子代数，则当$A$同时也是$H$的子余代数，并且反元$S$将$A$映射到$A$中时，称$A$为一个霍普夫子代数。换言之，当把$H$的乘法、余乘法、余元和反元限制到$A$上时，$A$本身就是一个霍普夫代数（此外还要求$H$的单位元 $1$属于$A$）。Warren Nichols 与 Bettina Zoeller 于 1989 年证明的 **Nichols–Zoeller 自由性定理表明：如果$H$是有限维的，则自然的$A$-模 $H$ 是有限秩的自由模；这推广了子群的拉格朗日定理【9】。由此及积分理论推出，半单有限维霍普夫代数的霍普夫子代数自动是半单的。

在霍普夫代数$H$中，若一个霍普夫子代数$A$满足右稳定性条件，即$\operatorname{ad}_r(h)(A) \subseteq A \quad \text{对所有 } h \in H$，则称$A$在$H$中是右正规的。其中右伴随映射$\operatorname{ad}_r$定义为$\operatorname{ad}_r(h)(a) = S(h_{(1)})\, a \, h_{(2)}, \quad a \in A, \; h \in H$.类似地，若$A$在左伴随映射$\operatorname{ad}_l(h)(a) = h_{(1)}\, a \, S(h_{(2)})$下稳定，则称$A$在$H$中是左正规的。当反元$S$是双射时，左右正规条件等价，此时称$A$为一个正规霍普夫子代数。

在$H$中的正规霍普夫子代数$A$满足以下条件（作为$H$的子集相等）：$H A^{+} = A^{+} H$,其中$A^{+}$表示$A$上余元的核。此正规性条件意味着$H A^{+}$是$H$的一个霍普夫理想（即：它既是余元核中的代数理想，也是余代数的余理想，并且在反元下稳定）。因此，可以得到一个商霍普夫代数$H / H A^{+}$,以及一个满射$H \to H / A^{+} H$,其理论与群论中正规子群和商群的理论类似【10】。
\subsubsection{霍普夫整环}
设$R$是一个整环，$K$是其分式域。若$H$是$K$上的一个霍普夫代数，则$R$上的一个霍普夫整环$O$是$H$ 的一个整环，且在代数与余代数运算下封闭：特别地，余乘法$\Delta : O \to O \otimes O$成立【11】。
\subsubsection{类群元素}
一个类群元素是满足$\Delta(x) = x \otimes x$的非零元素 $x$。类群元素在反元给出的逆元运算下构成一个群【12】。

一个原始元素$x$ 则满足：$\Delta(x) = x \otimes 1 + 1 \otimes x$【13】【14】。 
\subsection{例子}
\begin{table}[ht]
\centering
\caption{请输入表格标题}\label{tab_HPFds1}
\begin{tabular}{|c|c|c|c|c|c|c|c|}
\hline
 & \textbf{取决于} & \textbf{余乘法} & \textbf{余元} & \textbf{反元} & \textbf{交换性} & \textbf{余交换性} & \textbf{备注} \\
\hline
群代数 $KG$ & 群 $G$ & $\Delta(g) = g \otimes g \quad \forall g \in G$ & $\varepsilon(g) = 1 \quad \forall g \in G$ & $S(g) = g^{-1} \quad \forall g \in G$ & 当且仅当 $G$ 是阿贝尔群时成立 & 是 &  \\
\hline
从有限群到 $K$ 的函数 $f$，$K^G$（带点加与点乘运算） & 有限群 $G$ & $\Delta(f)(x,y) = f(xy)$ & $\varepsilon(f) = f(1_G)$ & $S(f)(x) = f(x^{-1})$ & 是 & 当且仅当 $G$ 是阿贝尔群时成立 &  \\
\hline
紧群上的表象函数 & 紧群 $G$ & $\Delta(f)(x,y) = f(xy)$ & $\varepsilon(f) = f(1_G)$ & $S(f)(x) = f(x^{-1})$ & 是 & 当且仅当 $G$ 是阿贝尔群时成立 & 反之，每一个在 $\mathbb{C}$ 上具有有限 Haar 积分的交换的、对合的、约化的霍普夫代数都以这种方式出现，这给出了Tannaka–Krein 对偶性的一种表述【15】。 \\
\hline
代数群上的正则函数 &  & $\Delta(f)(x,y) = f(xy)$ & $\varepsilon(f) = f(1_G)$ & $S(f)(x) = f(x^{-1})$ & 是 & 当且仅当 $G$ 是阿贝尔群时成立 & 反之，域上的每一个交换霍普夫代数都以这种方式来自一个群概形，从而给出了范畴之间的一种反等价【16】。 \\
\hline
张量代数 $T(V)$ & 向量空间 $V$ & $\Delta(x) = x \otimes 1 + 1 \otimes x, \; x \in V,\; \Delta(1) = 1 \otimes 1$ & $\varepsilon(x) = 0$ & $S(x) = -x \quad \forall x \in T^1(V)$（并推广到更高阶张量幂） & 当且仅当 $\dim(V) = 0,1$ 时成立 & 是 & 对称代数与外代数（它们是张量代数的商）在此余乘法、余元和反元定义下也都是霍普夫代数 \\
\hline
普遍包络代数 $U(\mathfrak{g})$ & 李代数 $\mathfrak{g}$ & $\Delta(x) = x \otimes 1 + 1 \otimes x \quad \forall x \in \mathfrak{g}$（此规则与对易子相容，因此可以唯一地扩展到整个 $U$） & $\varepsilon(x) = 0 \quad \forall x \in \mathfrak{g}$（同样可扩展到 $U$） & $S(x) = -x$ & 当且仅当 $\mathfrak{g}$ 是阿贝尔李代数时成立 & 是 &  \\
\hline
Sweedler 的霍普夫代数$H = K[c, x] / (c^2 = 1, \; x^2 = 0, \; xc = -cx)$。 & $K$ 是一个特征不等于 2 的域。 & $\Delta(c) = c \otimes c,\;\; \Delta(x) = c \otimes x + x \otimes 1,\;\; \Delta(1) = 1 \otimes 1$ & $\varepsilon(c) = 1,\;\; \varepsilon(x) = 0$ & $S(c) = c^{-1} = c,\;\; S(x) = -cx$ & 否 & 否 & 其底层向量空间由 $\{1, c, x, cx\}$ 张成，因此维数为 4。这是最小的同时非交换且非余交换的霍普夫代数例子。 \\
\hline
对称函数环【17】 &  & 以完全齐次对称函数 $h_k \; (k \geq 1)$ 表示$\Delta(h_k) = 1 \otimes h_k + h_1 \otimes h_{k-1} + \cdots + h_{k-1} \otimes h_1 + h_k \otimes 1$ & $\varepsilon(h_k) = 0$ & $S(h_k) = (-1)^k e_k$ & 是 & 是 &  \\
\hline
\end{tabular}
\end{table}
注意：有限群上的函数可以与群环对应起来，尽管它们更自然地被视为对偶结构——群环由有限个元素的和组成，因此可以通过将函数作用在这些求和元素上，与群上的函数配对。
\subsubsection{李群的上同调}
李群$G$的上同调代数（在域$K$上）是一个霍普夫代数：乘法由杯积给出；余乘法则由群乘法$G \times G \to G$所诱导：
$$
H^{*}(G,K) \;\longrightarrow\; H^{*}(G \times G,K) \;\cong\; H^{*}(G,K) \otimes H^{*}(G,K) .~
$$
这一观察实际上正是霍普夫代数概念的来源。利用这种结构，霍普夫证明了关于李群上同调代数的一个结构定理。

\textbf{定理（Hopf）}【18】设$A$是一个有限维、分次交换、分次余交换的霍普夫代数，定义在特征为 0 的域上。那么，作为代数，$A$是一个自由外代数，其生成元均为奇次数。
\subsection{表示论}
设$A$是一个霍普夫代数，$M, N$是$A$-模。则$M \otimes N$也是一个$A$-模，定义为：
$$
a (m \otimes n) := \Delta(a)(m \otimes n) = (a_{1} \otimes a_{2})(m \otimes n) = (a_{1}m \otimes a_{2}n),~
$$
其中$m \in M, \; n \in N,\; \Delta(a) = (a_1, a_2)$。此外，可以将平凡表示定义为基域 $K$，其作用为：
$$
a(m) := \epsilon(a) m, \quad m \in K.~
$$
最后，还可以定义$A$的对偶表示：若$M$是一个$A$-模，且 $M^{*}$ 是其对偶空间，则有：
$$
(af)(m) := f(S(a)m), ~
$$
其中 $f \in M^{*}, \; m \in M$。

$\Delta, \epsilon, S$ 三者之间的关系保证了某些自然的向量空间同态实际上是 $A$-模同态。例如：向量空间的自然同构$M \to M \otimes K$与$M \to K \otimes M$ 同时也是 $A$-模同构；向量空间同态$M^{*} \otimes M \to K,\; f \otimes m \mapsto f(m)$也是 $A$-模同态；但映射$M \otimes M^{*} \to K$并不一定是$A$-模同态。
\subsection{相关概念}
分次霍普夫代数常用于代数拓扑：它们是$H$-空间所有同调或上同调群的直和上自然出现的代数结构。

局部紧量子群是霍普夫代数的推广，并携带拓扑结构。李群上所有连续函数所成的代数就是一个局部紧量子群。

拟霍普夫代数是霍普夫代数的推广，其中余结合律只在“扭曲”意义下成立。它们被用于研究Knizhnik–Zamolodchikov 方程【21】。

乘子霍普夫代数由 Alfons Van Daele 于 1994 年提出【22】，是霍普夫代数的推广：其余乘法是从一个代数（可含幺元，也可无幺元）映射到该代数与自身张量积的乘子代数。

霍普夫群-(余)代数由 V. G. Turaev 于 2000 年提出，也是霍普夫代数的一种推广。
\subsubsection{弱霍普夫代数}
弱霍普夫代数，又称为量子群胚，是霍普夫代数的推广。与霍普夫代数一样，弱霍普夫代数构成一个自对偶的代数类；即：若$H$是一个（弱）霍普夫代数，则其对偶空间$H^*$（即$H$上的线性形式空间）也是一个（弱）霍普夫代数，其中的代数–余代数结构由$H$与其余代数–代数结构的自然配对给出。
\begin{itemize}
\item 通常，弱霍普夫代数$H$被假设为一个有限维的代数和余代数，带有余积$\Delta: H \to H \otimes H$和余元$\epsilon: H \to k$，它满足霍普夫代数的所有公理，除了可能存在 $\Delta(1) \neq 1 \otimes 1$ 或 $\epsilon(ab) \neq \epsilon(a)\epsilon(b)$（某些 $a,b \in H$）。取而代之，要求满足以下条件：
$$
(\Delta(1)\otimes 1)(1\otimes \Delta(1)) \;=\; (1\otimes \Delta(1))(\Delta(1)\otimes 1) \;=\; (\Delta \otimes \mathrm{Id})\Delta(1),~
$$
$$
\epsilon(abc) \;=\; \sum \epsilon(ab_{(1)}) \, \epsilon(b_{(2)}c) \;=\; \sum \epsilon(ab_{(2)}) \, \epsilon(b_{(1)}c),~
$$
对所有$a,b,c \in H$成立。
\item 此外，$H$拥有一个弱化的反元$S: H \to H$，它满足以下公理：\\
1.$S(a_{(1)})a_{(2)} = 1_{(1)} \, \epsilon(a1_{(2)}), \quad \forall a \in H$（右侧是一个投影，通常记作 $\Pi_R(a)$ 或 $\epsilon_s(a)$，其像是一个称为 $H_R$ 或 $H_s$ 的可分子代数）；\\
2.$a_{(1)} S(a_{(2)}) = \epsilon(1_{(1)}a) \, 1_{(2)}, \quad \forall a \in H$（这是另一个投影，通常记作 $\Pi_L(a)$ 或 $\epsilon_t(a)$，其像是一个可分代数 $H_L$ 或 $H_t$，并且通过 $S$ 与 $H_L$ 反同构）；\\
3.$S(a_{(1)})a_{(2)}S(a_{(3)}) = S(a), \quad \forall a \in H$.需要注意的是，如果 $\Delta(1) = 1 \otimes 1$，那么这些条件便化简为霍普夫代数反元的通常两条条件。
\end{itemize}
这些公理部分上是为了保证$H$-模范畴是一个刚性单oidal范畴。其中的单位 $H$-模就是上文提到的可分代数 $H_L$。

例如，有限群胚代数就是一个弱霍普夫代数。特别地，集合$[n]$上的群胚代数，若在其中每对元素$i, j \in [n]$之间都有一对可逆箭头$e_{ij}$与$e_{ji}$，则它同构于$n \times n$矩阵代数$H$。在这个特殊的$H$上，其弱霍普夫代数结构由下式给出：余积：$\Delta(e_{ij}) = e_{ij} \otimes e_{ij}$，余元：$\epsilon(e_{ij}) = 1$，反元：$S(e_{ij}) = e_{ji}$。
在此情形下，可分子代数 $H_L$ 与 $H_R$ 重合，并且它们是非中心的交换代数（即对角矩阵的子代数）。

关于弱霍普夫代数的早期理论性贡献，可以参见【23】和【24】。
\subsubsection{霍普夫代数胚}
\begin{itemize}
\item 参见 Hopf algebroid。
\end{itemize}
\subsubsection{与群的类比}
群可以通过与霍普夫代数相同的图表（或等价地，相同的运算）来公理化，只是此时 $G$被视为一个集合而非模。在这种情况下：
\begin{itemize}
\item 域$K$被替换为一个单点集；
\item 存在一个自然的余元（映射到单点）；
\item 存在一个自然的余乘法（对角映射）；
\item 单位元是群的单位元；
\item 乘法就是群的乘法；
\item 反元就是群的逆元。
\end{itemize}
从这个角度来看，可以把群看作是定义在“一元域”上的一个霍普夫代数【25】。
\subsection{编织单oidal范畴中的霍普夫代数}
霍普夫代数的定义可以自然推广到任意编织单oidal范畴【26】【27】。在这样一个范畴$(C, \otimes, I, \alpha, \lambda, \rho, \gamma)$中，一个霍普夫代数是一个六元组$(H, \nabla, \eta, \Delta, \varepsilon, S)$，其中 $H$是$C$中的一个对象，并且：\\
$\nabla : H \otimes H \to H$ —— 乘法；\\
$\eta : I \to H$ —— 单位；\\
$\Delta : H \to H \otimes H$ —— 余乘法；\\
$\varepsilon : H \to I$ —— 余元；\\
$S : H \to H$ —— 反元；

这些都是范畴$C$中的态射，并满足相应的条件。

1）三元组$(H, \nabla, \eta)$是单oidal范畴$(C, \otimes, I,\alpha, \lambda, \rho, \gamma)$中的一个幺半群，即下列图表交换【b】：
\begin{figure}[ht]
\centering
\includegraphics[width=14.25cm]{./figures/738aea6d6792a468.png}
\caption{} \label{fig_HPFds_2}
\end{figure}
2）三元组$(H, \Delta, \varepsilon)$是单oidal范畴$(C, \otimes, I, \alpha, \lambda, \rho, \gamma)$中的一个余幺半群，即下列图表交换【b】：
\begin{figure}[ht]
\centering
\includegraphics[width=14.25cm]{./figures/97f7c0702b2d47cd.png}
\caption{} \label{fig_HPFds_3}
\end{figure}
3）在$H$上的幺半群与余幺半群结构是相容的：乘法$\nabla$与单位$\eta$是余幺半群的态射，并且（在这种情形下是等价的）余乘法$\Delta$与余元 $\varepsilon$ 同时也是幺半群的态射；这意味着下列图表必须交换：
\begin{figure}[ht]
\centering
\includegraphics[width=14.25cm]{./figures/82afc9f712255c6b.png}
\caption{} \label{fig_HPFds_4}
\end{figure}
